% from ega4-9-fr.tex
\oldpage[IV-3]{93}

la fonction $x\mapsto\dim_x(W_{f(x)})$ est constructible en vertu de
(9.9.1, (i)), elle ne prend qu'un nombre fini de valeurs, donc les nombres
$\dim(W_{f(x)})$ ont une borne supérieure finie $m$ lorsque $x$ parcourt $X$~;
comme
$\operatorname{coprof}((\mathcal F_{f(x)})_x)\leq
\dim((\mathcal F_{f(x)})_x)\leq\dim(W_{f(x)})$, on voit que l'on a
$Z_n=\varnothing$ pour $n\geq m$. Enfin, il résulte de (6.11.2, (i)) que pour
tout $n$ et tout $s\in S$, $(Z_n)_s$ est \emph{fermé} dans $X_s$. D'après
(5.7.4), l'ensemble des $x\in X$ où $(\mathcal F_{f(x)})_x$ possède la
propriété $(S_k)$ est l'ensemble des $x\in X$ vérifiant toutes les relations
\begin{equation}
\tag{9.9.2.2}
\label{IV.9.9.2.2-fr}
\operatorname{codim}_x((Z_n)_{f(x)},W_{f(x)})>n+k
\end{equation}
pour tout $n\geq0$~; comme cette relation est automatiquement vérifiée pour
$n\geq m$, on n'a à considérer les relations (9.9.2.2) que pour $0\leq n<m$.
Mais en vertu de (9.9.1, (ii)), l'ensemble $V_{n,k}$ des $x$ vérifiant (9.9.2.2)
est constructible, et il en est de même de l'intersection de ces ensembles pour
$0\leq n<m$.

(ix) L'assertion résulte ici aussitôt de (vii), l'ensemble des $x\in X$ où
$(\mathcal F_{f(x)})_x$ est un module de Cohen-Macaulay étant défini par la
relation $\operatorname{coprof}((\mathcal F_{f(x)})_x)=0$.

\begin{corollary}[9.9.3]
\label{IV.9.9.3-fr}
Soient $f:X\to S$ un morphisme de présentation finie, $\mathcal F$ un
$\mathcal O_X$-Module de présentation finie, $\mathbf P$ l'une des propriétés
(i) à (ix) de (9.9.2). Alors l'ensemble des $s\in S$ tels que la propriété
$\mathbf P$ soit vraie en tous les points $x\in X_s$, est localement
constructible dans $S$.
\end{corollary}

En effet, son complémentaire est l'image par $f$ du complémentaire de
l'ensemble $E$ des points où $\mathbf P$ est vraie. Comme $E$ est localement
constructible dans $X$, il en est de même de $X-E$, donc $f(X-E)$ est
localement constructible dans $S$ en vertu du théorème de Chevalley (1.8.4).

\begin{proposition}[9.9.4]
\label{IV.9.9.4-fr}
Soit $f:X\to S$ un morphisme localement de présentation finie. L'ensemble des
$x\in X$ tels que $X_{f(x)}$ ait au point $x$ l'une (déterminée) des propriétés
suivantes~:
\begin{enumerate}
\item[(i)] être géométriquement régulier~;
\item[(ii)] posséder la propriété $(R_k)$ géométrique~;
\item[(iii)] être géométriquement normal~;
\item[(iv)] être géométriquement réduit (i.e. séparable)~;
\item[(v)] être géométriquement ponctuellement intègre~;
\end{enumerate}
est une partie localement constructible de $X$.
\end{proposition}

Pour les propriétés (iv) et (v), cela résulte de (9.9.2, (iv) et (v)) appliqué à
$\mathcal F=\mathcal O_X$. Pour les autres propriétés, compte tenu de (6.7.8),
on se ramène, comme au début de (9.9.2) au cas où $S$ est noethérien et intègre
et de point générique $\eta=f(x)$. En outre, en vertu du théorème de platitude
générique (6.9.1), on peut, en remplaçant $S$ par un voisinage ouvert de
$\eta$, supposer que $f$ est un morphisme \emph{plat}. Dire que $X_{f(y)}$ est
géométriquement régulier au point $y$ signifie alors que $f$ est
\emph{régulier} au point $y$ (6.8.1), et on sait que l'ensemble $L$ de ces
points est \emph{ouvert} dans $X$ (6.8.7), ce qui prouve la proposition dans le
cas (i).

Pour prouver le cas (ii), posons $F=X-L$, qui est \emph{fermé} dans $X$. Dire
que $X_s$

